% Options for packages loaded elsewhere
\PassOptionsToPackage{unicode}{hyperref}
\PassOptionsToPackage{hyphens}{url}
\PassOptionsToPackage{dvipsnames,svgnames,x11names}{xcolor}
%
\documentclass[
  11pt,
]{article}
\usepackage{amsmath,amssymb}
\usepackage{iftex}
\ifPDFTeX
  \usepackage[T1]{fontenc}
  \usepackage[utf8]{inputenc}
  \usepackage{textcomp} % provide euro and other symbols
\else % if luatex or xetex
  \usepackage{unicode-math} % this also loads fontspec
  \defaultfontfeatures{Scale=MatchLowercase}
  \defaultfontfeatures[\rmfamily]{Ligatures=TeX,Scale=1}
\fi
\usepackage{lmodern}
\ifPDFTeX\else
  % xetex/luatex font selection
\fi
% Use upquote if available, for straight quotes in verbatim environments
\IfFileExists{upquote.sty}{\usepackage{upquote}}{}
\IfFileExists{microtype.sty}{% use microtype if available
  \usepackage[]{microtype}
  \UseMicrotypeSet[protrusion]{basicmath} % disable protrusion for tt fonts
}{}
\makeatletter
\@ifundefined{KOMAClassName}{% if non-KOMA class
  \IfFileExists{parskip.sty}{%
    \usepackage{parskip}
  }{% else
    \setlength{\parindent}{0pt}
    \setlength{\parskip}{6pt plus 2pt minus 1pt}}
}{% if KOMA class
  \KOMAoptions{parskip=half}}
\makeatother
\usepackage{xcolor}
\usepackage[margin=2.2cm]{geometry}
\usepackage{longtable,booktabs,array}
\usepackage{calc} % for calculating minipage widths
% Correct order of tables after \paragraph or \subparagraph
\usepackage{etoolbox}
\makeatletter
\patchcmd\longtable{\par}{\if@noskipsec\mbox{}\fi\par}{}{}
\makeatother
% Allow footnotes in longtable head/foot
\IfFileExists{footnotehyper.sty}{\usepackage{footnotehyper}}{\usepackage{footnote}}
\makesavenoteenv{longtable}
\setlength{\emergencystretch}{3em} % prevent overfull lines
\providecommand{\tightlist}{%
  \setlength{\itemsep}{0pt}\setlength{\parskip}{0pt}}
\setcounter{secnumdepth}{-\maxdimen} % remove section numbering
\ifLuaTeX
  \usepackage{selnolig}  % disable illegal ligatures
\fi
\usepackage{bookmark}
\IfFileExists{xurl.sty}{\usepackage{xurl}}{} % add URL line breaks if available
\urlstyle{same}
\hypersetup{
  pdftitle={Right-isosceles-free subsets of rectangular boards},
  colorlinks=true,
  linkcolor={Maroon},
  filecolor={Maroon},
  citecolor={Blue},
  urlcolor={Blue},
  pdfcreator={LaTeX via pandoc}}

\title{Right-isosceles-free subsets of rectangular boards}
\author{}
\date{}

\begin{document}
\maketitle

\section{Right-Isosceles-Free Subsets of Rectangular Grids: Exact Values
for Narrow
Boards}\label{right-isosceles-free-subsets-of-rectangular-grids-exact-values-for-narrow-boards}

\subsection{Abstract}\label{abstract}

Let a(m,n) be the largest number of points of an m×n grid (m rows, n
columns of lattice points) that contain no three points forming a
non-degenerate right isosceles triangle. We prove a closed form for
two-row boards, a(2,n) = 2⌈n/2⌉ (Theorem 1), and give a short proof of
the one-row case. We compute a(m,n) exactly by constraint programming
for many small boards, including all m ≤ 4 with n ≤ 8 and m = 5 with n ≤
12, and for m = 3 with n ≤ 18. For three-row boards we prove a(3,n) =
n+2 for n ≥ 4 (computer-assisted upper bound; see the three-row-boards
project (paper/main.tex)). We did not find a table of these rectangular
values in OEIS, but the check was limited (see Section 5).

\subsection{1. Definitions}\label{definitions}

Points are lattice points (x, y) with 0 ≤ x \textless{} n and 0 ≤ y
\textless{} m. Three points form a right isosceles triangle if, for some
vertex V and the other two vertices P and Q, the vectors P−V and Q−V are
perpendicular and of equal nonzero length. A set is \emph{admissible} if
it contains no such triple. Write a(m,n) for the maximum size of an
admissible subset.

\subsection{2. Results}\label{results}

\subsubsection{Theorem 1 (two-row and one-row
boards)}\label{theorem-1-two-row-and-one-row-boards}

For all n ≥ 1:

\begin{enumerate}
\def\labelenumi{(\roman{enumi})}
\tightlist
\item
  a(1,n) = n,
\item
  a(2,n) = 2⌈n/2⌉.
\end{enumerate}

\textbf{Proof.}

\begin{enumerate}
\def\labelenumi{(\roman{enumi})}
\item
  All points of a one-row board are collinear, and a non-degenerate
  triangle has three non-collinear vertices. So no triple is forbidden,
  and a(1,n) = n.
\item
  \emph{Upper bound.} Split the columns into consecutive pairs \{1,2\},
  \{3,4\}, \ldots{} If n is odd, the last column is alone. Each pair of
  columns forms a 2×2 block of lattice points. Any three of the four
  corners of an axis-parallel unit square form a right isosceles
  triangle. So each 2×2 block contains at most two points of an
  admissible set. Each leftover single column (when n is odd) contains
  at most two points. Summing over the ⌈n/2⌉ blocks or columns gives
  \textbar S\textbar{} ≤ 2⌈n/2⌉.
\end{enumerate}

\emph{Lower bound.} Take S = \{(x,y) : x even, 0 ≤ x \textless{} n, y ∈
\{0,1\}\}. This has 2⌈n/2⌉ points. Suppose P, Q, R ∈ S form a right
isosceles triangle with right angle at P. Write u = Q − P = (a,b) and R
− P = (−b,a) (the rotation of u by 90°, up to sign; the other
orientation is handled the same way). All points of S have even
x-coordinate, so the x-differences a and −b are even. All points of S
have y ∈ \{0,1\}, so b and a lie in \{−1, 0, 1\}. Hence a and b are even
and lie in \{−1,0,1\}, so a = b = 0. Then Q = P, which is degenerate. So
S is admissible. ∎

\subsubsection{Computed values}\label{computed-values}

The table lists values a(m,n) with m rows and n columns, for n ≥ m (the
cases computed). Every entry shown was reported OPTIMAL by the solver,
so the upper and lower bounds met. Entries for the 3×19 and 3×20 boards,
which the solver did not close, are discussed separately below.

\begin{longtable}[]{@{}
  >{\raggedright\arraybackslash}p{(\columnwidth - 32\tabcolsep) * \real{0.0588}}
  >{\raggedright\arraybackslash}p{(\columnwidth - 32\tabcolsep) * \real{0.0588}}
  >{\raggedright\arraybackslash}p{(\columnwidth - 32\tabcolsep) * \real{0.0588}}
  >{\raggedright\arraybackslash}p{(\columnwidth - 32\tabcolsep) * \real{0.0588}}
  >{\raggedright\arraybackslash}p{(\columnwidth - 32\tabcolsep) * \real{0.0588}}
  >{\raggedright\arraybackslash}p{(\columnwidth - 32\tabcolsep) * \real{0.0588}}
  >{\raggedright\arraybackslash}p{(\columnwidth - 32\tabcolsep) * \real{0.0588}}
  >{\raggedright\arraybackslash}p{(\columnwidth - 32\tabcolsep) * \real{0.0588}}
  >{\raggedright\arraybackslash}p{(\columnwidth - 32\tabcolsep) * \real{0.0588}}
  >{\raggedright\arraybackslash}p{(\columnwidth - 32\tabcolsep) * \real{0.0588}}
  >{\raggedright\arraybackslash}p{(\columnwidth - 32\tabcolsep) * \real{0.0588}}
  >{\raggedright\arraybackslash}p{(\columnwidth - 32\tabcolsep) * \real{0.0588}}
  >{\raggedright\arraybackslash}p{(\columnwidth - 32\tabcolsep) * \real{0.0588}}
  >{\raggedright\arraybackslash}p{(\columnwidth - 32\tabcolsep) * \real{0.0588}}
  >{\raggedright\arraybackslash}p{(\columnwidth - 32\tabcolsep) * \real{0.0588}}
  >{\raggedright\arraybackslash}p{(\columnwidth - 32\tabcolsep) * \real{0.0588}}
  >{\raggedright\arraybackslash}p{(\columnwidth - 32\tabcolsep) * \real{0.0588}}@{}}
\toprule\noalign{}
\begin{minipage}[b]{\linewidth}\raggedright
m ~n
\end{minipage} & \begin{minipage}[b]{\linewidth}\raggedright
1
\end{minipage} & \begin{minipage}[b]{\linewidth}\raggedright
2
\end{minipage} & \begin{minipage}[b]{\linewidth}\raggedright
3
\end{minipage} & \begin{minipage}[b]{\linewidth}\raggedright
4
\end{minipage} & \begin{minipage}[b]{\linewidth}\raggedright
5
\end{minipage} & \begin{minipage}[b]{\linewidth}\raggedright
6
\end{minipage} & \begin{minipage}[b]{\linewidth}\raggedright
7
\end{minipage} & \begin{minipage}[b]{\linewidth}\raggedright
8
\end{minipage} & \begin{minipage}[b]{\linewidth}\raggedright
9
\end{minipage} & \begin{minipage}[b]{\linewidth}\raggedright
10
\end{minipage} & \begin{minipage}[b]{\linewidth}\raggedright
11
\end{minipage} & \begin{minipage}[b]{\linewidth}\raggedright
12
\end{minipage} & \begin{minipage}[b]{\linewidth}\raggedright
13
\end{minipage} & \begin{minipage}[b]{\linewidth}\raggedright
14
\end{minipage} & \begin{minipage}[b]{\linewidth}\raggedright
15
\end{minipage} & \begin{minipage}[b]{\linewidth}\raggedright
16
\end{minipage} \\
\midrule\noalign{}
\endhead
\bottomrule\noalign{}
\endlastfoot
1 & 1 & 2 & 3 & 4 & 5 & 6 & 7 & 8 & & & & & & & & \\
2 & & 2 & 4 & 4 & 6 & 6 & 8 & 8 & 10 & 10 & 12 & 12 & 14 & 14 & 16 &
16 \\
3 & & & 4 & 6 & 7 & 8 & 9 & 10 & 11 & 12 & 13 & 14 & 15 & 16 & 17 &
18 \\
4 & & & & 6 & 8 & 9 & 10 & 12 & 13 & 14 & 16 & 16 & 18 & 19 & & \\
5 & & & & & 9 & 10 & 12 & 13 & 15 & 16 & 17 & 18 & & & & \\
\end{longtable}

The m = 3 row continues: a(3,19) ≥ 21 and a(3,20) ≥ 22 were found but
not proved optimal (the solver's upper bounds were 37 and 39). These are
consistent with the conjecture below.

\subsubsection{Observations and
conjecture}\label{observations-and-conjecture}

\textbf{Observation 2.} For 3×n boards with n ≥ 4 the computed values
satisfy a(3,n) = n+2 for every n from 4 to 18.

\textbf{Proposition 3 (computer-checked construction).} For n ≡ 2 (mod
4), with n ≥ 6, the set

S\_n = \{(x,y) : y ∈ \{0,2\}, x mod 4 ∈ \{0,1\}\}

is admissible in the 3×n board and has n+2 points. We checked this by
brute force for n = 6, 10, \ldots, 30. We did not write a general proof.

\textbf{Theorem 4 (proved, see the three-row-boards paper
(\texttt{paper/main.tex})).} For all n ≥ 4, a(3,n) = n+2. The upper
bound is computer-assisted (a verified potential certificate on 443
states); the lower bound is an explicit construction valid for all n.

Both bounds for all n are now proved in the three-row-boards paper
(\texttt{paper/main.tex}).

\textbf{Remark on the 2-row pattern.} Restricting 3×n sets to rows 0 and
2 reaches n+2 only when n ≡ 2 (mod 4). For n ≡ 0, 1, 3 (mod 4), the
optimal sets use the middle row, so no single two-row pattern gives the
general lower bound.

\subsection{3. Method}\label{method}

We model the problem as a constraint satisfaction problem. Each lattice
point is a Boolean variable. For each right isosceles triple, we add the
clause ``not all three points are chosen.'' We add the constraint that
each 2×2 block holds at most two chosen points, which is valid by the
argument in Theorem 1 and tightens the linear relaxation. We maximize
the number of chosen points with Google OR-Tools CP-SAT on 2 CPU cores,
with time limits from 60 to 300 seconds.

Every reported solution was checked by an independent routine. That
routine enumerates all triples and tests the right-isosceles condition
directly, without using the solver's triangle list. All stored solutions
for m ≤ 4, n ≤ 8 passed, as did the constructions in Theorem 1 and
Proposition 3.

\subsection{4. What is and isn't
established}\label{what-is-and-isnt-established}

\begin{itemize}
\tightlist
\item
  \textbf{Proved:} Theorem 1 (both parts), with complete arguments.
\item
  \textbf{Computed exactly:} every entry in the table marked OPTIMAL in
  the solver output, i.e.~all of m ≤ 2 up to the displayed range, m = 3
  with n ≤ 18, m = 4 with n ≤ 14, and m = 5 with n ≤ 12. These are
  finite computations, not general proofs.
\item
  \textbf{Not established:} a closed form for m ≥ 4, and the values of
  the m = 4 and m = 5 rows beyond the table. The m = 3 closed form is
  now proved in the three-row-boards paper (\texttt{paper/main.tex}).
\item
  \textbf{Asymptotics:} we do not address growth in n.~The square-grid
  literature (Károlyi--Solymosi, arXiv:2607.22828, 2026) gives F(n) =
  Ω(n\^{}1.3) and F(n) ≪ n²/(log n)\^{}(1+c) for the square problem.
  Rectangular boards with fixed m are a different regime.
\end{itemize}

\subsection{5. Novelty}\label{novelty}

We searched OEIS for right-isosceles-triangle subsets of rectangular
grids. The only related entry we found was A271906, which covers square
grids n×n with n ≤ 11, and no rectangular table. We did not search the
broader combinatorics literature. Theorem 1 is elementary, and it is
quite possible that it appears in the literature on the square-grid
problem or on rectangular board problems, so we make no novelty claim
for it. The computed table is the most likely new contribution, but that
is contingent on a fuller literature check.

\subsection{6. Reproducibility}\label{reproducibility}

Code and data are in \texttt{./}:

\begin{itemize}
\tightlist
\item
  \texttt{rect.py}: triangle enumeration and CP-SAT model for m×n
  boards.
\item
  \texttt{rect2.py}: runs for the larger boards.
\item
  \texttt{rect\_4x8.json}, \texttt{rect\_more.json}: solver output.
\item
  \texttt{verify.py}: independent checker (square grids; the rectangular
  checks used the same triple test).
\end{itemize}

\end{document}
